% реферат по теме 'Аутентификация и доступ'
\documentclass{SibFU-docs}

% доп. пакеты
\usepackage{graphicx}
\usepackage{float}
\usepackage{hyperref}
\usepackage{caption}
\usepackage{subcaption}
\usepackage{amsmath}
% предотвартим ошибку с \Bbbk перед загрузкой пакета amssymb
\let\Bbbk\relax
\usepackage{amssymb}
\usepackage{booktabs}
\usepackage{tabularx}
\usepackage{longtable}
\usepackage{enumitem}
\usepackage{multirow}

% гиперссылки
\hypersetup{
    colorlinks=true,
    linkcolor=black,
    citecolor=blue,
    urlcolor=blue
}

% библиография
\addbibresource{report.bib}

% титульный лист
\begin{document}
\setlist[itemize]{left=1.3cm, itemsep=1pt, topsep=2pt} % для выравнивания списков
\setlist[enumerate]{left=1.3cm, itemsep=1pt, topsep=2pt}
\renewcommand{\contentsname}{}
\mycovertitle
{Гуманитарный институт}
{Кафедра прикладной информатики в искусстве и интерактивном медиа}
{Аутентификация и доступ}
{ст. преподаватель И.~Р.~Нигматуллин}
{Логинова Ю.В, Захаров И.М., Пачковская А.В., Тимофеева А.А., Дегтярева В.В., Шикайкова К.С., Рябинин Т.В., Мельников С.В.}
 
\begin{center}
\bfseries РЕФЕРАТ
\end{center}
\noindent

Реферат по теме: «Аутентификация и доступ» содержит 32 страниц и 25 использованных источников.

Ключевые слова: АУТЕНТИФИКАЦИЯ, ИДЕНТИФИКАЦИЯ ПОЛЬЗОВАТЕЛЯ, УЧЁТНАЯ ЗАПИСЬ, УПРАВЛЕНИЕ ПАРОЛЯМИ, КОНТРОЛЬ ДОСТУПА, ПРАВА ПОЛЬЗОВАТЕЛЕЙ, РОЛЕВЫЕ МОДЕЛИ БЕЗОПАСНОСТИ, ДИСКРЕЦИОННЫЙ ДОСТУП, МАНДАТНЫЙ ДОСТУП, ЖУРНАЛИРОВАНИЕ ДЕЙСТВИЙ, МНОГОФАКТОРНАЯ АУТЕНТИФИКАЦИЯ, ИНФОРМАЦИОННАЯ БЕЗОПАСНОСТЬ

В работе исследуются фундаментальные принципы аутентификации и управления доступом в информационных системах. Рассмотрены базовые понятия учётной записи и структуры системы идентификации пользователя, включая процессы создания, верификации и управления идентификаторами. Особое внимание уделено методам аутентификации и управления паролями, анализу моделей безопасности — дискреционной, мандатной и ролевой, а также уровням прав доступа в операционных и сетевых системах. Проанализированы механизмы контроля доступа и журналирования действий пользователей, а также современные методы многофакторной аутентификации и их применение в реальных сценариях.

Основные выводы:

\begin{enumerate}
    \item Учётная запись пользователя, объединяющая уникальные идентификаторы (логин, UID) и метаданные прав доступа, является фундаментальной структурной единицей любой системы идентификации, обеспечивая связь между субъектом и объектами защиты, а также возможность аудита и контроля действий.
    
    \item Эффективность аутентификации напрямую зависит от применяемых механизмов защиты паролей. Хеширование с солью (bcrypt, Argon2) и политики сложности обеспечивают устойчивость к атакам перебором, тогда как многофакторная аутентификация (пароль + TOTP + биометрия) радикально снижает риски компрометации единственного фактора, доминируя в современных системах корпоративной безопасности.

    \item Разграничение прав доступа реализуется через модели безопасности: дискреционная (DAC) предоставляет гибкость владельцу ресурса, мандатная (MAC) обеспечивает строгую иерархию доверия, ролевая (RBAC) оптимальна для крупных организаций, где права определяются служебными функциями и ролями.

    \item Будущее систем безопасности связано с переходом к адаптивной и контекстно-зависимой аутентификации. Технологии анализа поведения (UBA), биометрического распознавания и машинного обучения создают динамическую модель рисков, кардинально меняя подход к защите в облачных средах, IoT и распределённых системах.

\end{enumerate}

Рекомендации:

\begin{itemize}[label=-]
    \item Для критически важных систем использовать многофакторную аутентификацию (MFA) и современные алгоритмы хеширования паролей (Argon2, scrypt). Для локальных учётных записей внедрять политики ротации паролей и блокировку после неудачных попыток входа.

    \item При проектировании систем доступа начинать с ролевой модели (RBAC), определяя минимально необходимые привилегии для каждой роли, постепенно внедряя атрибутивный контроль доступа (ABAC) для более гибкого управления в динамических средах.

    \item Учитывать растущий тренд на адаптивную аутентификацию при разработке новых систем. Изучать основы поведенческого анализа и контекстных факторов (геолокация, устройство, время) для создания риск-ориентированных политик доступа.

    \item При выборе моделей безопасности всегда оценивать специфику системы и угрозовую модель, обеспечивая максимальную совместимость с существующими стандартами (OAuth 2.0, OpenID Connect) и оптимальное соотношение между уровнем защиты и удобством пользователей.

\end{itemize}

\newpage

\thispagestyle{empty}
\begin{center}
\bfseries СОДЕРЖАНИЕ
\end{center}
\vspace{-6pt}
\tableofcontents
\clearpage

% фантомные боли для секций, чтобы не было нумераций кроме тела реферата
\section*{}
\vspace*{-2.5em}
\begin{center}\textbf{ВВЕДЕНИЕ}\end{center}
\phantomsection
\addcontentsline{toc}{section}{Введение}
Современные информационные системы функционируют в условиях постоянных угроз безопасности и требуют эффективных механизмов защиты данных и контроля действий пользователей. Ключевую роль в обеспечении информационной безопасности играет система аутентификации и управления доступом, позволяющая разграничивать права пользователей, предотвращать несанкционированное вмешательство и обеспечивать целостность цифровых ресурсов. Развитие технологий идентификации и верификации личности стало фундаментом для построения безопасных вычислительных и сетевых инфраструктур, применяемых как в корпоративных, так и в публичных цифровых средах.

Современное состояние проблемы. Сегодня проблема обеспечения надёжной аутентификации и управления доступом является одной из наиболее актуальных в контексте цифровизации общества. Повсеместное использование интернет-сервисов, облачных платформ и сетевых приложений обусловливает необходимость многоуровневой защиты учётных записей. Наряду с традиционными методами — такими как парольная аутентификация — активно развиваются биометрические, поведенческие и многофакторные модели идентификации. При этом остаются нерешёнными вопросы балансировки между безопасностью, удобством и конфиденциальностью пользователя. Отдельное значение приобретают технологии журналирования и аудита действий, формирующие основу для анализа инцидентов и выявления нарушений доступа.

Актуальность данной работы определяется тем, что аутентификация пользователя и контроль прав доступа являются базовыми компонентами любой системы защиты информации — от локальных операционных систем до распределённых сетевых инфраструктур и облачных решений. Понимание принципов идентификации, разграничения прав и моделей безопасности необходимо каждому специалисту в области информационных технологий для построения надёжных архитектур и предотвращения цифровых угроз.

Цель работы: исследовать теоретические основы аутентификации и управления доступом, рассмотреть модели безопасности и уровни прав пользователей в современных системах, а также проанализировать перспективные методы защиты, включая многофакторные подходы.

Задачи исследования: раскрыть принципы аналого-цифрового преобразования звукового сигнала (дискретизация, квантование) и теоретические основы данного процесса; проанализировать методы и алгоритмы сжатия аудиоданных, провести сравнительную характеристику технологий сжатия с потерями и без потерь; систематизировать знания о ключевых аудиоформатах и кодеках, определить области их эффективного применения; исследовать инструментарий и базовые техники цифрового редактирования и монтажа аудио; рассмотреть современные технологии объёмного и пространственного звука (Dolby Atmos, Ambisonics и др.) и перспективы их развития.

Методы исследования: раскрыть понятие учётной записи и описать структуру системы идентификации пользователей; рассмотреть принципы аутентификации и надёжного управления паролями; проанализировать роли, категории и уровни прав доступа в операционных и сетевых системах; систематизировать основные модели безопасности — дискреционную, мандатную и ролевую; изучить методы контроля доступа и журналирования действий пользователей; рассмотреть современные подходы к многофакторной аутентификации и их практическое применение.

\newpage
% вручную увеличить номер секции и записать в оглавление с номером,
% чтобы в теле не печатался автоматический номер, но в tableofcontents он остался
\refstepcounter{section}
\addcontentsline{toc}{section}{\protect\numberline{\thesection} Понятие учетной записи и структура системы идентификации пользователя}
\vspace*{-2.5em}
\begin{center}\textbf{\thesection\ ПОНЯТИЕ УЧЕТНОЙ ЗАПИСИ И СТРУКТУРА СИСТЕМЫ ИДЕНТИФИКАЦИИ ПОЛЬЗОВАТЕЛЯ}\end{center}
\vspace*{-1.5em}
\begin{center}\textbf{\thesection\ ПОНЯТИЕ УЧЕТНОЙ ЗАПИСИ И СТРУКТУРА СИСТЕМЫ ИДЕНТИФИКАЦИИ ПОЛЬЗОВАТЕЛЯ}\end{center}
\vspace*{-1.5em}
\ManualSubsection{Учетная запись как фундамент систем управления доступом}

Современные информационные системы, от операционных систем до облачных платформ, опираются на концепцию учётной записи в качестве основного механизма цифрового представления субъекта. Учётная запись (аккаунт) представляет собой структурированный набор атрибутов и параметров, который позволяет системе однозначно идентифицировать пользователя, хранить его профильные данные, а также управлять его правами и полномочиями \cite{pro32_2025_uchetnaya_zapis}. Создание такой записи является отправной точкой для любого взаимодействия пользователя с защищёнными ресурсами, будь то регистрация на веб-сервисе или добавление сотрудника в корпоративный домен. Ключевые функции учётной записи простираются далеко за рамки простой идентификации. Как отмечают эксперты, она служит центральным узлом для персонализации среды, контроля доступа в соответствии с политиками безопасности, ведения журналов аудита для отслеживания действий и распределения системных ресурсов \cite{identity_blitz_2025_policies}. Таким образом, учётная запись формирует цифровую тень пользователя, связывая его личность с набором разрешённых операций в киберпространстве.

\ManualSubsection{Структура и компоненты учётной записи}

Архитектурно учётная запись состоит из нескольких взаимосвязанных компонентов, каждый из которых выполняет строго определённую роль. Базовым элементом является уникальный идентификатор (логин), например, имя пользователя или адрес электронной почты, который система использует для распознавания субъекта в процессе идентификации. Вторым критически важным компонентом выступают данные для аутентификации, чаще всего — хэш пароля, который служит криптографическим доказательством подлинности предъявленного идентификатора \cite{pro32_2025_uchetnaya_zapis}. Третьим структурным блоком являются атрибуты, определяющие права и полномочия. Сюда входят как прямые назначения привилегий, так и ссылки на роли или группы в системах ролевого (RBAC) или дискреционного (DAC) управления доступом. Наконец, учётная запись содержит метаданные: временные отметки создания и последнего использования, статус активности (активна/заблокирована), контактную информацию и историю действий. Особое внимание в контексте безопасности уделяется привилегированным и сервисным учётным записям, правам которых требуют жёсткого контроля и непрерывного мониторинга из-за высоких рисков их компрометации \cite{keepersecurity_2025_service_accounts}.

\ManualSubsection{Идентификация, аутентификация и авторизация: триединый процесс управления доступом}

Взаимодействие пользователя с системой регулируется последовательностью трёх фундаментальных процессов: идентификации, аутентификации и авторизации. Несмотря на кажущуюся схожесть терминов, каждый процесс решает свою уникальную задачу в цепочке предоставления доступа \cite{yandex_practicum_2025_ident_auth}. \textbf{Идентификация} — это заявление субъекта о себе путём предоставления уникального идентификатора (логина). На этом этапе система лишь проверяет существование учётной записи с таким идентификатором, отвечая на вопрос «Кто вы?». Процесс может быть первичным (при первой регистрации) или вторичным (при каждом последующем входе) \cite{rtsolar_2025_identifikacija}. \textbf{Аутентификация} — процедура проверки заявленной идентичности, ответ на вопрос «Докажите, что это вы?». Пользователь предоставляет доказательство, обычно секретный фактор (пароль, одноразовый код, биометрический образец), которое сверяется с эталоном, хранящимся в учётной записи. Именно здесь реализуется принцип многофакторности (MFA), требующий подтверждения из нескольких категорий: знания (пароль), владения (токен) и свойства (биометрия) \cite{rtsolar_2025_avtorizacija}. \textbf{Авторизация} — заключительный этап, на котором система определяет, какие именно действия и ресурсы разрешены аутентифицированному пользователю. Это ответ на вопрос «Что вам разрешено делать?». Решение принимается на основе прав, связанных с учётной записью, через механизмы ролевых или атрибутных моделей доступа (RBAC, ABAC) \cite{yandex_practicum_2025_ident_auth}. Чёткое разделение этих этапов является краеугольным камнем для построения безопасных и управляемых систем контроля доступа.

\newpage
\ManualSubsection{Система управления идентификацией и жизненный цикл учётной записи}

Для управления тысячами учётных записей в корпоративной среде используются специализированные системы управления идентификацией и доступом (Identity and Access Management, IAM). Архитектура такой системы включает хранилища данных (каталоги вроде Active Directory или LDAP), службы аутентификации, механизмы авторизации и средства аудита. Основной задачей IAM-системы является автоматизация всего жизненного цикла учётной записи: от её создания (provisioning) и назначения прав в соответствии с должностью сотрудника, через управление изменениями (например, при смене роли), до своевременной блокировки и удаления (deprovisioning) при увольнении \cite{identity_blitz_2025_policies}. Соблюдение формальных политик управления учётными записями является критически важным для безопасности. Эти политики регламентируют требования к сложности и периодичности смены паролей, порядок регулярного пересмотра назначенных прав (recertification), необходимость соблюдения принципа наименьших привилегий, а также обеспечивают своевременную деактивацию неиспользуемых или избыточных аккаунтов \cite{keepersecurity_2025_service_accounts}. Таким образом, учётная запись, будучи базовым кирпичиком цифровой идентичности, в рамках комплексной IAM-системы становится управляемым активом, безопасность которого поддерживается на всех этапах её существования.

\newpage
\refstepcounter{section}
\addcontentsline{toc}{section}{\protect\numberline{\thesection}Основы аутентификации и управления паролями}
\begin{center}\textbf{\thesection\ ОСНОВЫ АУТЕНТИФИКАЦИИ И УПРАВЛЕНИЯ ПАРОЛЯМИ}\end{center}
\vspace*{-1.5em}
\ManualSubsection{Фундаментальная триада AAA: идентификация, аутентификация, авторизация}

Базовым концептуальным каркасом для любого механизма контроля доступа является триада «AAA», включающая три последовательных и логически связанных процесса: идентификацию, аутентификацию и авторизацию \cite{keeper_auth_vs_authz_2023}. \textbf{Идентификация} представляет собой первичное заявление субъекта о своей личности, например, путём предоставления уникального логина или имени пользователя. На данном этапе система лишь осуществляет поиск соответствующей записи в своей базе данных, отвечая на вопрос «Кто вы?». Следующим обязательным шагом выступает \textbf{аутентификация} — процедура верификации предъявленной идентичности. Пользователь должен предоставить одно или несколько доказательств (факторов), подтверждающих, что он является владельцем заявленной учётной записи. Только после успешного прохождения аутентификации система переходит к этапу \textbf{авторизации}, на котором определяется совокупность конкретных прав и разрешений, предоставляемых пользователю для работы с ресурсами \cite{keeper_auth_vs_authz_2023}. Таким образом, идентификация устанавливает личность, аутентификация её доказывает, а авторизация определяет границы дозволенного.

\ManualSubsection{Механизмы аутентификации: от паролей к многофакторности}

Наиболее исторически распространённым методом аутентификации остаётся использование паролей, относящихся к категории фактора знания (something you know). Однако уязвимость паролей к множеству атак, таких как перебор по словарю, фишинг или использование утёкших учётных данных (credential stuffing), обусловила необходимость введения дополнительных защитных барьеров. Современным стандартом является \textbf{многофакторная аутентификация (MFA)}, требующая от пользователя предъявления двух или более независимых факторов из различных категорий: знания (пароль), владения (аппаратный токен, мобильное устройство для получения одноразового кода) или свойства (биометрические данные) \cite{nist_passwords_2024}. Эволюция аутентификации движется в сторону \textbf{беспарольных методов}, среди которых ключевую роль играют стандарты FIDO2 и WebAuthn. Данные технологии основаны на криптографии с открытым ключом: на устройстве пользователя создаётся уникальная пара ключей, а для входа используется биометрическая проверка или PIN-код, что обеспечивает высочайший уровень защиты от фишинга и перехвата учётных данных.

\ManualSubsection{Безопасное хранение и криптографическая защита паролей}

Принципиальным требованием безопасности на стороне сервера является запрет на хранение паролей пользователей в открытом (cleartext) виде. Вместо этого системы хранят криптографические \textbf{хэши} паролей — результат необратимого математического преобразования. Простого хэширования недостаточно из-за угрозы атак с использованием предвычисленных радужных таблиц. Для нейтрализации этой угрозы применяется \textbf{соль (salt)} — уникальная случайная строка, генерируемая для каждого пароля и добавляемая к нему перед хэшированием, что гарантирует уникальность хэша даже для идентичных паролей. Для противодействия аппаратному перебору используются специально разработанные \textbf{медленные алгоритмы хэширования (KDF)}, такие как Argon2, bcrypt или PBKDF2. Эти алгоритмы сознательно требуют значительных вычислительных ресурсов и времени, что делает массовый перебор паролей экономически и практически нецелесообразным, практически не затрудняя при этом легитимные операции входа.

\ManualSubsection{Современные стандарты и рекомендации по управлению паролями}

Современные подходы к политикам парольной безопасности, отражённые в авторитетных руководствах, таких как специальная публикация NIST SP 800-63B, претерпели значительную эволюцию \cite{nist_passwords_2024}. Ключевым трендом стал отход от жёстких, но зачастую контрпродуктивных правил. Вместо требования частой обязательной смены пароля (например, каждые 90 дней), что ведёт к созданию пользователями предсказуемых последовательностей, акцент смещён на контроль длины и проверку паролей по спискам известных уязвимых комбинаций \cite{nist_passwords_2024}. Длина парольной фразы признаётся более важным фактором криптостойкости, чем сложность из-за включения специальных символов. Настоятельно рекомендуется поощрять пользователей к применению \textbf{менеджеров паролей}, которые позволяют генерировать, безопасно хранить и использовать уникальные сложные пароли для каждого сервиса, полностью решая проблему их повторного использования и запоминания. В корпоративной среде управление доступом всё чаще делегируется централизованным системам на основе протоколов \textbf{OAuth 2.0} и \textbf{OpenID Connect (OIDC)}, которые позволяют предоставлять доступ к приложениям без передачи им самим паролей от основной учётной записи, минимизируя поверхность атаки и риски при компрометации одного из сервисов.

\newpage
\refstepcounter{section}
\addcontentsline{toc}{section}{\protect\numberline{\thesection}Роли и уровни прав доступа в операционных и сетевых системах}
\begin{center}\textbf{\thesection\ РОЛИ И УРОВНИ ПРАВ ДОСТУПА В ОПЕРАЦИОННЫХ И СЕТЕВЫХ СИСТЕМАХ}\end{center}
\vspace*{-1.5em}
\begin{center}\textbf{\thesection\ РОЛИ И УРОВНИ ПРАВ ДОСТУПА В ОПЕРАЦИОННЫХ И СЕТЕВЫХ СИСТЕМАХ}\end{center}
\vspace*{-1.5em}

\ManualSubsection{Роль как абстракция для управления доступом}

В операционных и сетевых системах роль представляет собой именованный набор разрешений и полномочий, предназначенный для назначения пользователям в соответствии с их функциональными обязанностями. Абстракция роли является ключевой для модели ролевого управления доступом (RBAC), так как отделяет управление правами от управления индивидуальными идентификаторами. Основное преимущество ролевого подхода заключается в значительном упрощении администрирования: вместо настройки сотен разрешений для каждого пользователя администратор определяет роли один раз, после чего назначение пользователя на роль автоматически наделяет его всеми связанными привилегиями. Это также обеспечивает соответствие организационной структуре, позволяет легко управлять изменениями при смене должности сотрудника и снижает вероятность ошибок, связанных с избыточным или недостаточным назначением прав \cite{microsoft2025rbac}. Один пользователь может одновременно обладать несколькими ролями, что позволяет гибко комбинировать полномочия в соответствии с его обязанностями.

\ManualSubsection{Иерархия уровней привилегий}

Современные операционные системы реализуют многоуровневую архитектуру привилегий, которая обеспечивает изоляцию между критическими компонентами системы и пользовательскими приложениями. На аппаратном уровне это выражается в концепции колец защиты, где наиболее привилегированное кольцо 0 зарезервировано для ядра операционной системы. На логическом уровне в пользовательском пространстве устанавливается иерархия от наименьших к наибольшим правам: гостевая учётная запись, обычный пользователь, привилегированный пользователь и, наконец, администратор или суперпользователь (root). Фундаментальный принцип, лежащий в основе этой иерархии, — принцип наименьших привилегий, согласно которому любой субъект должен обладать минимальным набором прав, необходимым для выполнения его легитимных задач \cite{cycode2024polp}. Следование этому принципу минимизирует потенциальный ущерб в случае компрометации учётной записи или выполнения вредоносного кода.

\ManualSubsection{Модель прав доступа в Unix-подобных системах}

Классическая модель управления доступом в Unix-подобных системах основана на триаде «владелец — группа — остальные» для каждого файла или каталога. Каждому элементу этой триады могут быть назначены три базовых типа разрешений: чтение (r), запись (w) и выполнение (x). Когда пользователь пытается получить доступ к объекту, система проверяет его соответствие категориям в строгой последовательности: сначала проверяются права владельца, затем — права группы, и только в последнюю очередь — права для всех остальных \cite{nersc2012unix}. Эта простая, но эффективная модель позволяет гибко управлять доступом в многопользовательской среде: владелец сохраняет полный контроль над своими данными, группа обеспечивает совместную работу над проектами, а права для остальных определяют публичный доступ. Для централизованного управления пользователями и группами часто используются службы каталогов, такие как LDAP \cite{harvard2025permissions}.

\ManualSubsection{Управление доступом в экосистеме Windows}

Windows использует более сложную и детализированную модель управления доступом, основанную на списках управления доступом (ACL). Права назначаются не по трём фиксированным категориям, а индивидуально для пользователей и групп безопасности. В автономных системах используются локальные группы (например, «Администраторы», «Пользователи»), в то время как в корпоративных средах с Active Directory управление централизуется через доменные группы. Ключевым механизмом распространения политик и членства в группах являются объекты групповой политики (GPO), которые позволяют автоматически добавлять, например, членов группы «Администраторы домена» в локальную группу администраторов на всех компьютерах домена \cite{woshub2025localadmin}. Эта модель обеспечивает высокую степень гранулярности и централизации, необходимую для управления крупными гетерогенными инфраструктурами.

\ManualSubsection{Взаимодействие сетевых и локальных разрешений}

При организации сетевого доступа к файловым ресурсам в Windows критически важным является понимание взаимодействия двух систем разрешений: разрешений общего доступа и разрешений файловой системы NTFS. Разрешения общего доступа (Share Permissions) регулируют доступ к ресурсу по сети, в то время как разрешения NTFS (NTFS Permissions) управляют доступом к файлам и папкам на локальном диске. Когда пользователь обращается к ресурсу по сети, Windows применяет наиболее ограничивающие права из двух наборов: если через общий доступ разрешено «чтение», а через NTFS — «изменение», результатом будет только право на чтение \cite{netwrix2025ntfs}. Рекомендуемой практикой является установка максимально разрешающих прав на общий доступ (например, «Полный контроль» для группы «Все») и тонкая настройка безопасности с помощью разрешений NTFS, что упрощает администрирование и минимизирует риск ошибок.

\ManualSubsection{Принцип разделения обязанностей}

Принцип разделения обязанностей является важным дополнением к ролевому управлению доступом и принципу наименьших привилегий. Его суть заключается в том, что ни один пользователь не должен обладать полным набором привилегий, достаточным для злоупотребления критически важной бизнес-функцией \cite{botha2004separation}. Для этого потенциально конфликтующие полномочия распределяются между несколькими независимыми ролями или пользователями. Например, в финансовой системе роль, создающая платёжные поручения, и роль, утверждающая их, должны быть назначены разным сотрудникам. Реализация этого принципа может быть статической (запрет на одновременное назначение конфликтующих ролей) или динамической (проверка на конфликт в контексте выполнения конкретной операции), что предотвращает мошенничество и ошибки, снижая операционные риски.

\ManualSubsection{Механизмы временного повышения привилегий}

Следование принципу наименьших привилегий в повседневной работе обеспечивается механизмами временного повышения прав, которые позволяют выполнять отдельные привилегированные операции без постоянной работы под учётной записью с максимальными правами. В Unix-подобных системах для этого используется утилита `sudo`, которая позволяет авторизованным пользователям выполнять команды от имени суперпользователя или другого пользователя, что обеспечивает детальный аудит всех действий. В Windows аналогичную функцию выполняет контроль учётных записей пользователей, который запрашивает подтверждение или учётные данные администратора при попытке внесения изменений, затрагивающих систему. Эти механизмы создают дополнительный барьер для несанкционированных действий и вредоносного программного обеспечения, минимизируя окно уязвимости и способствуя соблюдению политик безопасности.

\newpage
\refstepcounter{section}
\addcontentsline{toc}{section}{\protect\numberline{\thesection}Модели безопасности: дискреционная, мандатная и ролевая}
\begin{center}\textbf{\thesection\ МОДЕЛИ БЕЗОПАСНОСТИ: ДИСКРЕЦИОННАЯ, МАНДАТНАЯ И РОЛЕВАЯ}\end{center}
\vspace*{-1.5em}
\ManualSubsection{Концепция моделей управления доступом}

Модель управления доступом представляет собой формализованный набор принципов и механизмов, определяющих порядок предоставления субъектам (пользователям, процессам) прав на выполнение операций с объектами (файлами, данными, сервисами) в информационной системе. Она является ядром политики безопасности, связывая процессы аутентификации (проверки подлинности) и авторизации (предоставления прав). В современных условиях выбор адекватной модели является стратегической задачей, от которой зависит как уровень защищённости активов, так и эффективность бизнес-процессов организации \cite{rt_solar_models_2023}. Основные модели — дискреционная, мандатная и ролевая — представляют собой эволюционный ответ на растущие требования к безопасности и управляемости в системах различного масштаба и назначения.

\ManualSubsection{Дискреционная модель управления доступом}

Дискреционная модель управления доступом основана на принципе передачи полномочий по управлению доступом владельцу информационного ресурса. Владелец (чаще всего создатель файла или папки) обладает свободой («дискрецией») самостоятельно решать, кому и какие права (чтение, запись, выполнение) предоставить. Технически эта модель часто реализуется через списки управления доступом, которые привязываются к каждому объекту \cite{cleverence_comparison_2025}. Основным достоинством DAC является её гибкость и простота понимания, что делает её удобной для небольших рабочих групп и персонального использования. Ключевым недостатком выступает децентрализация контроля: решения владельца могут быть ошибочными или противоречить корпоративной политике, что создаёт риски несанкционированного распространения информации и сложности при централизованном аудите \cite{falcongaze_audit_2026}. Таким образом, DAC подходит для сред, где важна оперативность и сотрудничество, но не обеспечивает достаточного уровня гарантий для защиты критически важных данных.

\ManualSubsection{Мандатная модель управления доступом}

Мандатная модель является полной противоположностью дискреционной с точки зрения централизации и строгости. В этой модели ни один пользователь, включая владельца ресурса, не может самостоятельно изменять правила доступа. Права жёстко регламентируются централизованной политикой безопасности, определяемой администраторами в соответствии с формальными требованиями (например, государственными стандартами). Основу MAC составляют метки безопасности, которые присваиваются как пользователям (уровень допуска), так и объектам данных (степень секретности). Доступ разрешается только в том случае, если уровень допуска субъекта соответствует или превышает метку объекта, что реализует принципы «не читать вверх» и «не писать вниз» для предотвращения утечки информации \cite{rt_solar_models_2023}. Такая модель, несмотря на сложность администрирования и ограниченную гибкость, является незаменимой в средах с жёсткими требованиями к конфиденциальности, таких как государственные или военные системы, где человеческий фактор в принятии решений о доступе должен быть полностью исключён.

\ManualSubsection{Ролевая модель управления доступом}

Ролевая модель управления доступом была разработана как ответ на потребности крупного корпоративного сектора в балансе между безопасностью, управляемостью и масштабируемостью. Её ключевая идея заключается в введении промежуточного абстрактного звена — роли. Права доступа назначаются не конкретным пользователям, а ролям, которые соответствуют должностным функциям в организации («бухгалтер», «менеджер по закупкам»). Пользователи получают права путём назначения на одну или несколько ролей. Такой подход кардинально упрощает администрирование прав в больших организациях: при смене должности сотрудника достаточно изменить его роль, а не перенастраивать сотни индивидуальных разрешений \cite{cleverence_comparison_2025}. RBAC естественным образом поддерживает принцип наименьших привилегий, обеспечивая пользователей только теми правами, которые необходимы для выполнения их рабочих задач, что ограничивает потенциальный ущерб от компрометации учётной записи. Благодаря своей структурированности и относительной простоте управления, ролевая модель стала фактическим стандартом для корпоративных информационных систем и широко поддерживается современными системами управления идентификацией \cite{rt_solar_models_2023}.

\ManualSubsection{Сравнительный анализ и выбор модели}

Выбор модели управления доступом является критически важным архитектурным решением, которое должно основываться на анализе требований конкретной организации. Для объективного сравнения можно выделить ключевые критерии: степень централизации контроля, гибкость управления, сложность администрирования и уровень гарантий безопасности. Дискреционная модель предлагает максимальную гибкость, но минимальные гарантии из-за децентрализованного управления. Мандатная модель обеспечивает максимальный уровень безопасности за счёт жёсткой централизации, но является наименее гибкой и наиболее сложной в администрировании. Ролевая модель занимает сбалансированное промежуточное положение, сочетая приемлемый уровень безопасности с высокой управляемостью в условиях большого числа пользователей \cite{cleverence_comparison_2025}.

На практике в сложных системах часто применяется гибридный подход. Например, в операционной системе может использоваться ролевая модель для управления системными функциями, дискреционная — для пользовательских файлов в домашних каталогах, а доступ к корпоративным данным в облачной среде может регулироваться более современными атрибутными политиками. Рекомендуется начинать проектирование с ролевой модели как базовой для корпоративного сегмента, дополняя её элементами других моделей там, где это диктуется специфическими требованиями безопасности или бизнес-процессов \cite{falcongaze_audit_2026}. Понимание сильных и слабых сторон каждой модели позволяет проектировать адекватные, эффективные и поддерживаемые системы защиты информации.

\newpage
\refstepcounter{section}
\addcontentsline{toc}{section}{\protect\numberline{\thesection}Контроль доступа и основы журналирования действий пользователей}
\begin{center}\textbf{\thesection\ КОНТРОЛЬ ДОСТУПА И ОСНОВЫ ЖУРНАЛИРОВАНИЯ ДЕЙСТВИЙ ПОЛЬЗОВАТЕЛЕЙ}\end{center}
\vspace*{-1.5em}
\ManualSubsection{Фундаментальные цели контроля доступа и аудита}

В архитектуре информационной безопасности системы контроля доступа и журналирования (аудита) выполняют взаимодополняющие проактивную и реактивную функции. Контроль доступа представляет собой проактивный механизм, предназначенный для защиты конфиденциальности и целостности данных путём предотвращения несанкционированных действий. Его работа основана на реализации классической триады «AAA»: идентификации, аутентификации и авторизации пользователя. Напротив, журналирование является реактивным инструментом, обеспечивающим подотчётность всех операций в системе. Его основная цель заключается в создании неизменяемой хронологической записи событий для последующего расследования инцидентов, анализа производительности и доказательства соответствия регуляторным требованиям, таким как GDPR, PCI DSS или стандартам ФСТЭК \cite{Stallings2018, Peltier2016}. В совокупности эти механизмы образуют замкнутый цикл безопасности, где политики доступа предотвращают нарушения, а аудит позволяет обнаружить их обход и усовершенствовать защиту.

\ManualSubsection{Эволюция моделей контроля доступа}

Развитие моделей контроля доступа отражает поиск баланса между требованиями безопасности, гибкостью управления и сложностью администрирования в системах разного масштаба и назначения. Дискреционная модель, исторически первая, предоставляет владельцу ресурса свободу назначать права другим пользователям, что обеспечивает простоту, но несёт риски из-за децентрализованного и потенциально неконтролируемого распространения прав \cite{Tanenbaum2006}. Мандатная модель, применяемая в системах с жёсткими требованиями к защите государственной тайны, полностью централизует управление через метки безопасности, исключая человеческий фактор, но жертвуя гибкостью. Ролевая модель, ставшая корпоративным стандартом, вводит абстракцию роли для упрощения управления правами тысяч пользователей, естественным образом поддерживая принцип наименьших привилегий. Наиболее современной и гибкой является атрибутная модель, которая принимает решения о доступе динамически, на основе анализа множества характеристик пользователя, ресурса, действия и контекста, что особенно востребовано в сложных облачных и распределённых средах \cite{Stallings2018}. Каждая модель занимает свою нишу, а в реальных системах часто применяются их гибридные комбинации.

\ManualSubsection{Принципы и практика системного журналирования}

Эффективная система журналирования должна строиться на чётких принципах, гарантирующих полноту, достоверность и сохранность записей. Критически важными для логирования являются события аутентификации (все успешные и неудачные попытки входа), управления доступом (изменение учётных записей и прав), операции с конфиденциальными данными, а также все действия, совершённые привилегированными пользователями \cite{Bejtlich2013}. Для структурирования записей используется система уровней важности, такая как Syslog, которая позволяет разделять события от аварийных до отладочных. Ключевым организационным принципом является централизованный сбор логов с использованием специализированных агентов и хранилищ, что обеспечивает независимость аудиторского следа от скомпрометированного источника и упрощает анализ. Параллельно должны реализовываться меры защиты целостности (хеширование, цифровые подписи), конфиденциальности (шифрование) и неизменяемости записей, например, с помощью носителей с функцией однократной записи \cite{Peltier2016}.

\ManualSubsection{Интеграция и анализ: роль SIEM-систем}

Простого сбора и хранения логов недостаточно для противодействия современным сложным атакам. Системы класса SIEM решают задачу корреляции событий из множества разнородных источников: журналов операционных систем, приложений, сетевого оборудования и средств защиты. Аналитическое ядро SIEM, опираясь на предустановленные и настраиваемые правила, способно выявлять неочевидные многоэтапные атаки. Например, система может автоматически сгенерировать инцидент высокой важности, связав в единую цепочку события: множественные неудачные попытки подбора пароля с одного IP-адреса, последующую успешную аутентификацию с того же адреса и немедленный доступ с этой учётной записи к критически важным файлам \cite{Bejtlich2013}. Таким образом, SIEM превращает сырой поток событий в оперативную осведомлённость о безопасности, обеспечивая возможность быстрого реагирования на угрозы.

\ManualSubsection{Сопутствующие принципы безопасности}

Для создания целостного защищённого контура контроль доступа и аудит должны поддерживаться рядом фундаментальных организационных и технических принципов. Принцип наименьших привилегий требует, чтобы каждый субъект системы обладал минимальным набором прав, необходимым для выполнения его функций, что радикально ограничивает потенциальный ущерб от компрометации. Принцип разделения обязанностей, заимствованный из финансового контроля, предотвращает злоупотребления, распределяя критические операции между несколькими независимыми субъектами \cite{Peltier2016}. Технической основой для корректного аудита выступает синхронизация времени всех устройств инфраструктуры по протоколу NTP, без которой невозможна точная корреляция событий. Наконечь, жизненный цикл безопасности завершается принципом регулярного пересмотра, который предписывает проводить периодические аудиты назначенных прав и анализ политик для их постоянной актуализации в соответствии с изменяющимися бизнес-процессами и угрозами.

\newpage
\refstepcounter{section}
\addcontentsline{toc}{section}{\protect\numberline{\thesection}Современные методы многофакторной аутентификации и их применение}
\begin{center}\textbf{\thesection\ СОВРЕМЕННЫЕ МЕТОДЫ МНОГОФАКТОРНОЙ АУТЕНТИФИКАЦИИ И ИХ ПРИМЕНЕИЕ}\end{center}
\vspace*{-1.5em}
\ManualSubsection{Фундаментальные принципы многофакторной аутентификации}

Многофакторная аутентификация представляет собой метод контроля доступа, требующий от пользователя представления двух или более независимых доказательств своей подлинности из различных категорий факторов. Классическая модель опирается на три фундаментальных категории: фактор знания, которым владеет только пользователь, фактор владения, основанный на обладании физическим или программным устройством, и фактор сущности, использующий уникальные биометрические характеристики. Ключевым принципом MFA является комбинация факторов из разных категорий, так как использование двух факторов из одной категории не обеспечивает необходимого повышения уровня безопасности \cite{yandex_cloud_2fa}. Данный подход, рекомендованный стандартами NIST, радикально снижает вероятность успеха несанкционированного доступа, так как для компрометации аккаунта злоумышленник должен преодолеть несколько независимых защитных барьеров одновременно.

\ManualSubsection{Технологии второго фактора: от OTP к push-уведомлениям}

Исторически первым и наиболее распространённым методом второго фактора стали одноразовые пароли, доставляемые через SMS. Однако данный метод обладает существенными уязвимостями, такими как риск перехвата, атаки SIM-swapping и зависимость от операторов связи, что побуждает к переходу на более безопасные альтернативы \cite{anti_malware_mfa_2025}. Более надёжным решением являются одноразовые пароли на основе времени, генерируемые локально на устройстве пользователя по алгоритму TOTP. TOTP-коды, создаваемые приложениями-аутентификаторами, не требуют сетевого соединения для получения и устойчивы к перехвату при передаче. Дальнейшим развитием в сторону удобства стали push-уведомления, которые позволяют подтвердить вход одним нажатием в мобильном приложении, однако они создают риск атак типа «MFA Fatigue», при которых пользователя засыпают запросами в надежде на случайное подтверждение.

\ManualSubsection{Аппаратные токены и стандарты беспарольной аутентификации}

Высший уровень безопасности обеспечивают аппаратные токены, такие как YubiKey или отечественные Рутокены, работающие на основе криптографии с открытым ключом. Эти устройства физически хранят приватный ключ, который никогда не покидает токен, и подписывают им запрос на аутентификацию. Такая модель полностью защищает от фишинга, так как ключ криптографически привязан к конкретному домену сайта \cite{yandex_cloud_2fa}. Аппаратные ключи являются основой современных стандартов FIDO2 и WebAuthn, которые определяют будущее аутентификации, двигаясь в сторону полностью беспарольных сценариев. Стандарт WebAuthn, поддерживаемый всеми основными браузерами, позволяет веб-приложениям интегрировать безопасную аутентификацию с помощью биометрии или аппаратных ключей, минимизируя зависимость от традиционных паролей.

\ManualSubsection{Биометрия и адаптивные системы}

Биометрическая аутентификация, использующая уникальные физиологические или поведенческие характеристики, обеспечивает высокий уровень удобства и безопасности. Современные методы включают распознавание отпечатков пальцев, лица и голоса, причём системы всё чаще используют биометрию в качестве основного или единственного фактора, храня на сервере не сами биометрические образцы, а их криптографические шаблоны \cite{anti_malware_mfa_2025}. Эволюционным шагом стало внедрение адаптивной аутентификации, в которой система динамически оценивает уровень риска каждой попытки входа на основе контекстных атрибутов: местоположения пользователя, устройства, времени и поведенческих паттернов. При низком уровне риска доступ может быть предоставлен после минимальной проверки, тогда как подозрительная активность активирует требование дополнительных факторов. Эта модель, часто усиленная алгоритмами машинного обучения, обеспечивает оптимальный баланс между безопасностью и удобством пользователя.

\ManualSubsection{Корпоративное внедрение и регуляторные требования}

Внедрение MFA в корпоративной среде существенно отличается от личного использования и предполагает решение задач централизованного управления, интеграции с существующей инфраструктурой и соблюдения регуляторных норм. Корпоративные решения обеспечивают администраторам единую панель для управления политиками, мониторинга статуса пользователей, принудительного применения методов аутентификации в зависимости от роли сотрудника и генерации отчётов для аудита \cite{itglobal_mfa_corporate}. Критически важной является интеграция MFA с системами каталогов, такими как Active Directory, и поддержка протоколов единого входа для обеспечения беспрепятственного доступа к множеству корпоративных приложений. Стимулом для внедрения выступают не только внутренние политики безопасности, но и внешние регуляторные требования, такие как PCI DSS, ФСТЭК, HIPAA или GDPR, которые прямо предписывают использование MFA для защиты критически важных данных и систем.

\ManualSubsection{Тенденции развития и новые угрозы}

Эволюция MFA движется по пути отказа от уязвимых методов в пользу более безопасных и удобных, стандартизации на основе открытых протоколов и интеграции с системами анализа рисков. Ключевым трендом является переход к полностью беспарольным экосистемам на базе стандартов FIDO2 и концепции Passkeys, которые устраняют самый слабый элемент — человеческий фактор в создании и хранении паролей \cite{anti_malware_mfa_2025}. Параллельно развиваются и методы атак: от относительно простого фишинга и подмены SIM-карт до изощрённых атак на основе социальной инженерии и утомления пользователя. Это обуславливает необходимость постоянного пересмотра политик безопасности, обучения пользователей и использования комбинации методов, где аппаратные ключи и биометрия играют роль самого защищённого последнего рубежа, особенно для привилегированных учётных записей и доступа к критически важным активам.

\newpage
\section*{}
\vspace*{-2.5em}
\begin{center}\textbf{ЗАКЛЮЧЕНИЕ}\end{center}
\phantomsection
\addcontentsline{toc}{section}{Заключение}

Системы аутентификации и управления доступом являются ключевым элементом информационной безопасности и устойчивости современных цифровых инфраструктур. В ходе работы были рассмотрены основные принципы построения систем идентификации, аутентификации и разграничения прав доступа, а также сравнительный анализ моделей безопасности. Выявлено, что эффективность защиты определяется не только техническими средствами, но и комплексным применением организационных мер, политик и постоянного мониторинга активности пользователей.

Аутентификация выступает первой линией обороны, предотвращающей несанкционированное проникновение, тогда как управление доступом обеспечивает соблюдение принципов наименьших привилегий и структурированность взаимодействий между субъектами и объектами системы. Особую ценность представляет журналирование действий пользователей, позволяющее отслеживать и анализировать события безопасности в реальном времени. Это обеспечивает прозрачность процессов и способствует оперативному реагированию на инциденты.

Современные тенденции безопасности направлены на интеграцию интеллектуальных методов анализа, биометрических технологий и адаптивных алгоритмов принятия решений. Такие подходы позволяют не только усилить защищённость, но и повысить удобство работы конечных пользователей. Таким образом, система аутентификации и контроля доступа становится не статической технологией, а динамическим инструментом обеспечения доверия в цифровом пространстве.

\newpage
\begin{center}\textbf{СПИСОК СОКРАЩЕНИЙ}\end{center}
\phantomsection
\addcontentsline{toc}{section}{Список сокращений}

\vspace{1em}

{
\small
\noindent
\begin{tabular}{|l|p{0.75\textwidth}|}
\hline
\textbf{Сокращение} & \multicolumn{1}{c|}{\textbf{Расшифровка}} \\
\noalign{\hrule height 1.2pt}
\hline
MFA & Многофакторная аутентификация. Метод доступа, требующий минимум два независимых фактора подлинности.\\
\hline
2FA & Двухфакторная аутентификация. Частный случай MFA с использованием ровно двух факторов.\\
\hline
RBAC & Ролевое управление доступом. Модель, где права назначены ролям, а не пользователям.\\
\hline
DAC & Дискреционная модель управления доступом. Модель, где владелец ресурса сам решает, кому предоставить доступ.\\
\hline
MAC & Мандатная модель управления доступом. Строгая модель, управляемая централизованной политикой на основе меток секретности.\\
\hline
ABAC & Атрибутное управление доступом. Модель, где решение о доступе принимается на основе оценки динамических атрибутов.\\
\hline
ACL & Список управления доступом. Механизм, определяющий права пользователей и групп на операции с объектом.\\
\hline
IAM & Управление идентификацией и доступом. Комплекс процессов и технологий для управления цифровыми идентичностями и правами.\\
\hline
PoLP & Принцип наименьших привилегий. Правило предоставления субъекту минимально необходимых для работы прав.\\
\hline
SoD & Разделение обязанностей. Принцип распределения критических операций между разными лицами для предотвращения злоупотреблений.\\
\hline
OTP & Одноразовый пароль. Пароль, действительный для одной сессии аутентификации.\\
\hline
TOTP & Одноразовый пароль на основе времени. OTP, меняющийся через фиксированные временные интервалы (RFC 6238).\\
\hline
SIEM & Управление информацией и событиями безопасности. Система для сбора, анализа и отчётности по событиям безопасности.\\
\hline
LDAP & Облегчённый протокол доступа к каталогам. Протокол для доступа и управления централизованными службами каталогов.\\
\hline
OIDC & OpenID Connect. Протокол аутентификации поверх OAuth 2.0 для делегированной проверки подлинности.\\
\hline
\end{tabular}
}

\clearpage
\nocite{*}
\printbibliography[heading=bibliography]

\end{document}
